\chapter*{ВВЕДЕНИЕ}
\addcontentsline{toc}{chapter}{ВВЕДЕНИЕ}

Многопользовательские игровые приложения относятся к классу распределенных интерактивных систем, в которых несколько участников взаимодействуют с общим виртуальным состоянием через сеть. Корректность работы таких систем зависит не только от скорости передачи данных, но и от согласованности состояний, хранящихся и обрабатываемых на разных сторонах взаимодействия.

В условиях асинхронного сетевого обмена состояние игроков может рассогласовываться из-за задержек передачи сообщений, потери пакетов, дублирования данных, нарушения порядка доставки и различий во времени обработки событий. Для приложений, использующих UDP, эта проблема становится особенно существенной, так как данный протокол не предоставляет встроенных гарантий доставки и упорядочивания сообщений.

Одним из способов решения данной задачи является представление взаимодействия игроков в виде событийной модели с формально заданными состояниями и переходами. Использование конечных автоматов позволяет ограничить множество допустимых состояний системы, определить правила обработки событий и обеспечить одинаковый порядок изменения состояния у участников взаимодействия.

Целью выпускной квалификационной работы является разработка метода обеспечения согласованности состояний двух асинхронных игроков с помощью конечных автоматов.

Для достижения поставленной цели необходимо выполнить следующие задачи:
\begin{itemize}
	\item провести анализ предметной области синхронизации игровых состояний и передачи данных в многопользовательских играх;
	\item рассмотреть существующие подходы к синхронизации состояний и способы передачи игровых данных;
	\item описать типовые игровые данные, подлежащие хранению и передаче;
	\item определить основные причины рассинхронизации игровых состояний;
	\item разработать метод синхронизации состояний двух игроков с использованием конечных автоматов;
	\item определить формальную модель состояний системы и переходов между ними;
	\item разработать структуру конечного автомата, описывающего взаимодействие игроков;
	\item представить алгоритм синхронизации состояний в виде схем алгоритмов или диаграмм состояний;
	\item обосновать выбор технологий и средств разработки программной реализации метода;
	\item разработать программное обеспечение, реализующее предложенный метод синхронизации;
	\item разработать программное обеспечение, реализующее взаимодействие игроков;
	\item провести тестирование разработанного программного обеспечения;
	\item определить критерии оценки эффективности разработанного метода;
	\item провести исследование эффективности разработанного метода;
	\item оценить скорость синхронизации состояний и устойчивость алгоритма к различным сценариям взаимодействия игроков;
	\item сравнить результаты работы разработанного метода с существующими подходами синхронизации состояний.
\end{itemize}
